% !TeX root = ../papers/nc-p3.tex
\section{第二長方形の負垂直族の射影的な全商と二次錐}
\label{sec:zero-exchange-projective-quotient}
\begin{theorem}\label{thm:zero-exchange-projective-quotient}
$z=0,D=ad-bc\ne0$ の前節の負交換表示の全正則点では、
$[a:b:c:d]$ と $[d:b:c:a]$ だけが指定始点で同じ類を与える。
従って射影座標 $[p:b:c:u]$（$p=a+d,u=(a-d)^2$）で、全商は
$\mathbb P(1,1,1,2)$ の
\[
 p^2-u-4bc\ne0,\qquad (c-b,u-(p+2b)^2)\ne(0,0),\qquad
 (p+b,u-p^2)\ne(0,0)
\]
なる開集合である。周期1の軌跡は $p+c=0$ と $c=2b+p=0$ の和、それ以外の最小周期は2である。
反対階数は $16-[p=0]-2[c=0]-[p+2c=0]$ である。
この族の正則点と先の $z\ne0$ の負交換表示の間に指定始点の同型はない。
\end{theorem}
\begin{proof}[二つの射影変換と $c\ne0$]
命題\ref{prop:zero-exchange-vertical-identification}の四直線を
$a,d$ の交換で割ると記載の除外式になる。実際のテンソルで
\[
 I:[d:b:c:a],\qquad J:[-a-c:b+a+c+d:c:-d-c]
\]
がそれぞれ指定始点と一位置移動の変換である。各々で共通の四文字置換をした後
（位置ごとに別の置換を許す）、全非零係数の比は $\pm D^j$, $j=0,\pm1$ となる。
全整数文字関係を検査すると対角倍率が存在し、$D$ が単元なので
厚い係数の消滅を含めて恒等式を特殊化できる。
$c\ne0$ では第二長方形の $B=1$ の座標は $(-a/c,-d/c,b/c)$ である。
定理\ref{thm:parallel-vertical-orbits}の台代表への移動を $I,J$ で記録すれば、
指定始点の比較は $I$ だけで尽きる。代表の被覆は、先の座標を $(v,y,h)$ として
$v,y,1-v,1-y,h,v+y-1-h$ の零条件で確かめられる。
$I$ は $v,y$ を交換し、$J$ は $v,y$ を $1-v,1-y$ に、$h$ を $h+1-v-y$ に移す。
$v,y$ のうち0または1になる個数が0なら、後の二式の消滅数に従い147,145,142へ行く。
一個ならその値を $v=1$ にして139、または $h$ の特別値で134に行く
（$h=v+y-1$ は $vy=h$ により不可）。二個なら124,118,129のいずれかである。
ここで $v=1$ の規格化前に $J$ の適用を要した場合も、添字移動の回数を保存して
同定理の奇・偶両方の置換比較を使う。
\end{proof}

\begin{proof}[無限遠 $c=0$ の点と純曲線]
$ad\ne0$ である。左の点集合を前節の $P,Q$ を用いて調べる。
$a\ne d$ の場合、純一次元部分は
\[
 x_0^2=-\frac{b+d}{a(a-d)}x_2^2,\qquad
 x_1^2=-\frac{a+b}{d(a-d)}x_2^2
\]
の錐で、頂点は $e_3$、残る孤立点は $e_0\pm e_1$ の二点である。
二定数は同時に消えない。両方非零なら相異なる四直線、一方だけ零なら二本の二重直線となる。
実際、混合小行列式で $x_2x_3Q=x_0x_1P=0$ が得られる。
$Q\ne0$ で $x_2,x_3$ の一方が非零の場合を除く計算には、左側四例外直線に乗らないことを使う。
$Q=0,P\ne0$ のとき四つの混合積 $x_ix_jP$ ($i<2\le j$) は全て零である
（$z=0$ でも $m_{1235}$ が $-x_0x_3P$ を与える）。
前節の消去を辿ると、両式の外で可能なのは上の二孤立点だけである。
二式で定まるスキームに全小行列が消え、各直線の一般点では $x_2x_3\ne0$ から $Q$ が、
次に非零な $x_0,x_1$ の一方を含む混合式から $P$ が生成される。
従って埋込み点を除くと上の完全交叉そのものである。

頂点からの射影平面で四つの直線の方向が相異なれば、四方向の完全四辺形の三対角点が
座標点 $e_0,e_1,e_2$ である。孤立二点の像を結ぶ直線は第一・第二点だけを含む。
二重直線なら、方向スキームは二つの長さ2点である。その二接線の交点が消えた方の座標の頂点であり、
二支持の割線と孤立点の割線との交点がもう一方の低位の頂点となる。
二支持の割線上で後者の調和共役を取れば $e_2$ の像を得る。
元の空間で孤立点の割線に戻れば $e_0,e_1$ は交換を除いて決まる。$e_2$ の持上げだけには $e_3$ を加えられる。

$d=a$, $b\ne-a$ なら $P,Q$ は $x_0^2-x_1^2,x_2^2$ を張り、純曲線は二つの二重直線、
孤立点は存在しない。直線の一般点で $x_0x_1\ne0$ より $P$ を得て、$x_0x_3$ を含む式から
$Q$ を得る。台上で頂点以外に埋込みのある点は $e_0\pm e_1$ だけである。
前節の $P$ を法とする五式の除算でもこの両点上で $Q$ は生成されず、
両点と頂点以外では上の同じ二つの単元で $P,Q$ が生成される。
二直線の交点 $e_3$ と二つの埋込み点を枠に取る。残る $e_2$ の像を規定するには
節点枠の二座標 $v_0,v_1$ を使う。純理想は $(v_0v_1,x_2^2)$ であり、
この三点を保つ変換が $v_i$ に $x_2$ の倍を加えると積に二つの独立な項 $v_j x_2$ が生じる。
従って $e_2$ の持上げに許せるのはやはり $e_3$ の加法に限る。
奇位置では前記 $J$ の明示単項同型で同じ結果を使える。
四頂点の順の錐頂点は $(3,1,3,0)$、持上げを変えられる文字は $(2,0,2,1)$ である。
$a=d$ の場合の残りの二線は偶位置で $e_0\pm e_1$、奇位置で $e_2\pm e_3$ である。
直線の方向と頂点座標の零を保つ単項移動、または六式への直接代入で確かめられる。
\end{proof}

\begin{proof}[加法を許す枠の比較]
各位置の表示は、二つの固定文字の交換、文字倍率、持上げ文字に頂点の倍を加える四係数からなる。
四位置の頂点文字を一つ含む語が28個あり、二つ以上含む語はない。
従って持上げの変更は頂点語を固定し、頂点を含まない語に四列の行列の像を加える。
どちらの枠でも、全ての強制非零頂点語を許容台に送る交換は、恒等と四位置同時の交換だけである。
$a\ne d$ の場合には四列の台が元のテンソルの台と交わらず、指定始点の比較は
対角倍率と二つの置換に帰着する。36定数語と $a,b,d$ の各一語による格子差で、
$[a:b:d]$ が $I$ を除いて一意である。加法の四列は階数3で、その核が一次元加法群となる。
$a=d$ の場合は $a=1$ として、まず28語を固定する倍率のトーラスを取る。
そのSmith因子は $1^{11}$ で、残る零頂点語をこのトーラスの文字で分ける。
四列の像と残るテンソルは全て同じ原始文字の19語の空間にある。
四列が張る三次元を割った後、残るテンソルは $v+bw$ の形で $v,w$ は独立である。
従って射影化しても $b$ は一意である。同時交換は元の文字符号によって実現する。
この場合には一次元加法群のほかに一次元の射影対角安定群が残る。
\path{computations/verify_zero_exchange_projective_flags.py} が語の全置換、加法の列とその商を検算する。

$d=a,b=-a,c=0$ はただ一類である。$a=1$ とすれば前節の二次曲面表示を $z=0$ としたもので、
第一辺を対角化する行列はなお可逆、交代五平面の零化二ベクトルもなお非退化である。
従ってスキームとして二次錐 $x_0^2-x_1^2-x_2^2=0$ の対角グラフであり、全辺について $J$ で従う。
定理\ref{thm:quadric-cone-quotients}の余分な部分のない唯一の交代型に等しい。
これで滑らかな二次曲面枝の $z+1/z\to\infty$ における正則な特殊化も決まる。

$c\ne0$ の場合との重複は孤立点数で排除できる。$c\ne0$, $a\ne d$ なら
$P=Q=0$ の外に $e_0\pm e_1$ と $e_3$ の三点、$a=d$ なら $e_3$ だけが来る。
後者でも $\alpha=(a+b)(a+c),\ \beta=ac$ は同時に零でなく、二次式は完全交叉をなす。
小行列による外部の消去は上と同様である。
$z\ne0$ の表示とは三点・一点の孤立点数、または $c=0$ の場合の錐状曲線の交点・多重度や曲面の階数で区別できる。
商上の $J$ は $[p:b:c:u]\mapsto[-p-2c:b+p+c:c:u]$ だから、表示した周期1条件が従う。
反対階数は色0,1の二つずつの二行ブロックと、色2,3の行列式 $-4$ から得られる。
\end{proof}
